% Part: many-valued-logic
% Chapter: three-valued-logics
% Section: kleene

\documentclass[../../../include/open-logic-section]{subfiles}

\begin{document}

\olfileid{mvl}{thr}{skl}

\olsection{Logika-logika Kleene}

Stephen Kleene ngenalake rong logika telung nilai kang adhedhasar
pamikiran yen nilai kabeneran iku asil saka prosedur komputasi:
sawijining prosedur bisa ngasilake $\True$ utawa $\False$, nanging
uga bisa ora mandheg. Yen mangkono, nilai kabeneran kang cocog ora
katetepake, diwakili dening nilai~$\Undef$.

Kanggo ngitung negasi proposisi~$!A$, kita luwih dhisik ngitung
nilainé~$!A$, banjur ngasilake kosok baline. Yen komputasi $!A$
ora mandheg, sakabehe prosedur uga ora mandheg: mula negasi
$\Undef$ yaiku $\Undef$.

Kanggo ngitung konjungsi $!A \land !B$, ana rong cara. Cara kapisan,
kita bisa ngitung~$!A$ luwih dhisik, banjur $!B$; asile $\True$
yen asil loro-lorone~$\True$, lan $\False$ yen ora mangkono.
Yen salah siji komputasi ora mandheg, sakabehe prosedur uga ora
mandheg. Mula ing cara iki, yen salah siji komponen konjungsi
ora katetepake, konjungsine uga ora katetepake. Semono uga
kanggo disjungsi.

Nanging yen kita bisa ngetung $!A$ lan $!B$ bebarengan,
kita bisa oleh asil luwih cepet. Yen salah siji prosedur mandheg
lan ngasilake $\False$, kita bisa mandheg, amarga konjungsine
mesthi salah. Mula konjungsi kang nduweni siji komponen salah
iku salah, sanajan komponen liyane ora katetepake. Semono uga,
nalika ngitung disjungsi bebarengan, kita bisa mandheg sawise
salah siji prosedur ngasilake nilai bener: disjungsine mesthi
bener. Dadi ing tafsir iki kita bisa ngerti asil pratelan majemuk,
sanajan salah siji komponene ora katetepake. Ing kene $\Undef$
bisa diwaca ``durung dingerteni'', dudu ``ora katetepake''.

Rong tafsir iki ngasilake logika Kleene kang kuwat lan kang
ringkih. Kondisional ditetepake padha karo $\lnot !A \lor !B$.

\begin{defn}
\emph{Logika Kleene kang kuwat}~$\LogKs$ ditetepake nganggo matriks iki:
\begin{enumerate}
  \item Perangan tanpa tetapan saka basa proposisional baku
  $\Lang L_0$ kang ngemot $\lnot$, $\land$, $\lor$, $\lif$.
  \item Himpunan nilai kabeneran $V = \{\True, \Undef, \False\}$.
  \item $\True$ mung siji-sijine nilai pinilih, yaiku $V^+ = \{\True\}$.
  \item Fungsi kabeneran diwenehake dening tabel iki:
  \begin{center}
    \begin{tabular}{c|c} 
      $\tf{\lnot}$ & \\ 
      \hline  
      $\True$ & $\False$ \\ 
      $\Undef$ & $\Undef$ \\
      $\False$ & $\True$ 
    \end{tabular}
    \quad
    \begin{tabular}{c|ccc} 
      $\tf{\land}[\LogKs]$ & $\True$ & $\Undef$ & $\False$ \\ 
      \hline 
      $\True$ & $\True$ & $\Undef$ & $\False$ \\ 
      $\Undef$ & $\Undef$ & $\Undef$ & $\False$\\ 
      $\False$ & $\False$ & $\False$ & $\False$ 
    \end{tabular}
    \\[2ex]
    \begin{tabular}{c|ccc} 
      $\tf{\lor}[\LogKs]$ & $\True$ & $\Undef$ & $\False$ \\ 
      \hline 
      $\True$ & $\True$ & $\True$ & $\True$ \\ 
      $\Undef$ & $\True$ & $\Undef$ & $\Undef$ \\
      $\False$ & $\True$ & $\Undef$ & $\False$ 
    \end{tabular}
    \quad
    \begin{tabular}{c|ccc} 
      $\tf{\lif}[\LogKs]$ & $\True$ & $\Undef$ & $\False$ \\ 
      \hline 
      $\True$ & $\True$ & $\Undef$ & $\False$ \\ 
      $\Undef$ & $\True$ & $\Undef$ & $\Undef$  \\ 
      $\False$ & $\True$ & $\True$ & $\True$ 
    \end{tabular}
  \end{center} 
\end{enumerate}
\end{defn}

\begin{defn}
\emph{Logika Kleene kang ringkih}~$\LogKw$ ditetepake nganggo matriks iki:
\begin{enumerate}
  \item Perangan tanpa tetapan saka basa proposisional baku
  $\Lang L_0$ kang ngemot $\lnot$, $\land$, $\lor$, $\lif$.
  \item Himpunan nilai kabeneran $V = \{\True, \Undef, \False\}$.
  \item $\True$ mung siji-sijine nilai pinilih, yaiku $V^+ = \{\True\}$.
  \item Fungsi kabeneran diwenehake dening tabel iki:
  \begin{center}
    \begin{tabular}{c|c} 
      $\tf{\lnot}$ & \\ 
      \hline  
      $\True$ & $\False$ \\ 
      $\Undef$ & $\Undef$ \\
      $\False$ & $\True$ 
    \end{tabular}
    \quad
    \begin{tabular}{c|ccc} 
      $\tf{\land}[\LogKw]$ & $\True$ & $\Undef$ & $\False$ \\ 
      \hline 
      $\True$ & $\True$ & $\Undef$ & $\False$ \\ 
      $\Undef$ & $\Undef$ & $\Undef$ & $\Undef$\\ 
      $\False$ & $\False$ & $\Undef$ & $\False$ 
    \end{tabular}
    \\[2ex]
    \begin{tabular}{c|ccc} 
      $\tf{\lor}[\LogKw]$ & $\True$ & $\Undef$ & $\False$ \\ 
      \hline 
      $\True$ & $\True$ & $\Undef$ & $\True$ \\ 
      $\Undef$ & $\Undef$ & $\Undef$ & $\Undef$ \\
      $\False$ & $\True$ & $\Undef$ & $\False$ 
    \end{tabular}
    \quad
    \begin{tabular}{c|ccc} 
      $\tf{\lif}[\LogKw]$ & $\True$ & $\Undef$ & $\False$ \\ 
      \hline 
      $\True$ & $\True$ & $\Undef$ & $\False$ \\ 
      $\Undef$ & $\Undef$ & $\Undef$ & $\Undef$  \\ 
      $\False$ & $\True$ & $\Undef$ & $\True$ 
    \end{tabular}
  \end{center} 
\end{enumerate}
\end{defn}

\begin{prop}
  $\LogKs$ lan $\LogKw$ ora nduweni tautologi.
\end{prop}

\begin{proof}
  Yen $\pAssign v(p) = \Undef$ kanggo kabeh !!{propositional variable}s~$p$,
  saben formula~$!A$ ing perangan basa iki bakal nduweni nilai
  kabeneran~$\pValue v(!A) = \Undef$, amarga
  \[
    \tf{\lnot}(\Undef) = \tf{\lor}(\Undef, \Undef) = \tf{\land}(\Undef,
  \Undef) = \tf{\lif}(\Undef, \Undef) = \Undef
  \]
  ing loro logika. Amarga $\Undef \notin V^+$ ing $\LogKs$ lan
  $\LogKw$, miturut !!{valuation} iki $!A$ ora oleh nilai pinilih.
\end{proof}

Sanajan logika Kleene kang ringkih lan kang kuwat ora nduweni
tautologi, sesambungan konsekuensine ora trivial.

\begin{prob}
  Sesambungan endi ing ngisor iki kang lumaku ing (a) logika Kleene
  kang kuwat lan (b) logika Kleene kang ringkih? Wenehana tabel
  kabeneran kanggo saben sesambungan.
  \begin{enumerate}
    \item $p, p \lif q \Entails q$
    \item $p \lor q, \lnot p \Entails q$
    \item $p \land q \Entails p$
    \item $p \Entails p \land p$
    \item $p \Entails p \lor q$
  \end{enumerate}
\end{prob}

Dmitry Bochvar napsirake $\Undef$ minangka ``tanpa makna'' lan nyoba
nggunakake kanggo ngrampungake paradoks kayata paradoks Pembohong,
kanthi netepake ukara paradoks nduweni nilai~$\Undef$. Dheweke
ngenalake logika kang pokoke logika Kleene kang ringkih, nanging
ditambahi konektif liyane. Loro ing antarane yaiku ``negasi njaba''
lan operator ``ora katetepake'':
\begin{center}
  \begin{tabular}{c|c} 
    $\tf{\sim}$ & \\ 
    \hline  
    $\True$ & $\False$ \\ 
    $\Undef$ & $\True$ \\
    $\False$ & $\True$ 
  \end{tabular}
\quad
  \begin{tabular}{c|c} 
    $\tf{+}$ & \\ 
    \hline  
    $\True$ & $\False$ \\ 
    $\Undef$ & $\True$ \\
    $\False$ & $\False$ 
  \end{tabular}
\end{center}

\begin{prob}
  Apa $\sim$ ing logika Bochvar bisa ditetepake nganggo $\lnot$
  lan~$+$? Tegesé, apa ana !!a{formula} kang mung ngemot
  !!{propositional
  variable}~$p$ lan ora ngemot $\sim$, nanging
  mesthi nduweni nilai kabeneran padha karo~$\mathord{\sim}p$?
  Wenehana tabel kabeneran kanggo mbuktekake jawabanmu.
\end{prob}

\end{document}
